\section{Discussion}
\label{sec:disc}

\paragraph{What the evidence supports.}
We recommend that a benchmark report its visual headroom (blind ceiling, black-image accuracy and true-image accuracy), the accuracy on counter-cue questions where a cue is known, and, for a system, its accuracy with and without the image on the text-undecided stratum with the share of text-decisive questions.

\paragraph{Limits.}
\emph{Sufficiency.} The causal claim is for a planted shortcut, not for Med-CMR's own cue; the dose-response covers two benchmarks, a held-out prediction held for two further models, but the lock-out strength differs between benchmarks and is not explained by natural variation (Sec.~\ref{sec:general}), InternVL3-2B was not locked out at $0.99$, and larger models, other shortcuts and other tasks are untested.
\emph{Real shortcut.} Ten earlier attempts did not meet their criteria (Supp.~X), UL's criterion was chosen after them, and the cue is not the only barrier; the dependence that UL induces rests mostly on coarse image identity at 3B and less so at 7B.
\emph{Method.} The faithfulness audit uses reproductions of the organisers' judge at three sizes, not their judge (size unknown); no patient identifiers exist (case-disjointness and near-duplicate filtering only; thresholds were set by an AI assistant's review of montages); the released adapter's training membership is unknown; the analyses were implemented with an AI coding assistant under the hashed pre-registration.

\paragraph{Ethics and broader impact.}
We use MedReason~2026 (CC BY-NC 4.0, as stated by its organisers), PMC-VQA (CC BY-SA 4.0), the open-access subset of OmniMedVQA, A-OKVQA, MedXpertQA-MM~\cite{zuo2025medxpertqa} and MMMU~\cite{yue2024mmmu} under their published terms, for research only; nothing here supports clinical use.
Statements about released systems describe measured behaviour on stated benchmarks and inputs, not their clinical safety or their developers' intent, and the faithfulness audit uses reproductions of a judge, not the organisers' own.

\paragraph{Conclusion.}
\emph{Demonstrated:} in our controlled setting, a nonvisual shortcut locks the decision at a benchmark-dependent strength, and a penalty restores measurable image dependence for about 4 to 6 points of accuracy, partly on coarse image identity.

\emph{Not demonstrated:} a universal lock-out law, a mechanism explaining the difference between benchmarks, or restored semantic visual understanding.

\emph{Why it matters:} VQA evaluations should report visual headroom alongside accuracy, because high accuracy alone does not show that the image contributes.
